\documentclass[11pt,a4paper]{article}
\usepackage[margin=2.5cm]{geometry}
\usepackage{amsmath,amssymb,amsthm,mathtools,hyperref,enumitem}
\newtheorem{theorem}{Theorem}[section]
\newtheorem{corollary}[theorem]{Corollary}
\theoremstyle{definition}
\newtheorem{definition}[theorem]{Definition}
\newcommand{\R}{\mathbb{R}}
\newcommand{\E}{\mathbb{E}}
\newcommand{\Pbb}{\mathbb{P}}
\newcommand{\SCX}{SCX}
\newcommand{\rigorFull}{\textbf{[Rigor: Full Proof]}}
\newcommand{\Obs}{\mathcal{O}}

\title{SCX Causal Unidentifiability Theorem}
\author{SCX}
\date{2026}

\begin{document}
\maketitle

\begin{abstract}
Proves that causal direction claims are unidentifiable from observational data without declared identification assumptions. Three mutually incompatible Structural Causal Models produce identical observational distributions.
\end{abstract}
\subsection{Causal Counterfactual Corollary}
\label{subsec:causal_counterfactual}

The same identification barrier appears in ordinary causal claims.  A claim
such as ``$X$ causes $Y$'' is not a statement about the observed event alone:
it asserts a counterfactual contrast, e.g.\ that $Y$ would have been different
under an intervention $\mathrm{do}(X=x')$.  The counterfactual world is not
jointly observed with the factual world.  Hence the claim has empirical
content only after the analyst declares an identification strategy.

\begin{theorem}[SCX Causal Unidentifiability Theorem]
\label{thm:causal_unident}
Let $(X,Y,Z)$ take values in standard Borel spaces and let
$P_{XYZ}$ be any observed joint law admitting regular conditional
distributions.  Suppose the target causal claim is a counterfactual
statement about the effect of $X$ on $Y$, for example
$P(Y \in B \mid \mathrm{do}(X=x),Z=z)$ or an average treatment effect
$\psi = \mathbb{E}[Y_x - Y_{x'}]$.  Without a declared identification
assumption, the observational law $P_{XYZ}$ does not identify the causal
direction between $X$ and $Y$.  In particular, there exist structural causal
models inducing the same observed joint law $P_{XYZ}$ with the following
mutually incompatible causal interpretations:
\begin{enumerate}[label=(\alph*), leftmargin=*]
    \item $X \to Y$: the association is generated by a direct causal effect
    of $X$ on $Y$;
    \item $C \to X$ and $C \to Y$: the association is generated by a common
    cause $C$ whose observed component is $Z$;
    \item $Y \to X$: the association is generated by reverse causation.
\end{enumerate}
If $P_{Y\mid X,Z}$ or $P_{X\mid Y,Z}$ varies in the conditioning argument,
the corresponding directed model has an active controlled effect.  If the
observed association is only marginal, the same construction applies after
dropping $Z$ from the conditioning set and generating $Z$ from
$P_{Z\mid X,Y}$.
Consequently no statistic of observational samples from $P_{XYZ}$ can
verify which of these counterfactual interpretations is true.
\end{theorem}

\begin{proof}
\rigorFull
The proof is constructive.  Let $U_Z,U_X,U_Y$ be independent
$\mathrm{Unif}(0,1)$ random variables, and let
$Q_{A\mid B}$ denote a measurable conditional quantile or, more generally,
a randomization kernel that generates the conditional law $P(A\mid B)$.

\medskip\noindent
\textbf{World A: direct causation $X \to Y$.}
Define the SCM
\begin{align}
    Z_A &= Q_Z(U_Z), \nonumber\\
    X_A &= Q_{X\mid Z}(U_X \mid Z_A), \nonumber\\
    Y_A &= Q_{Y\mid X,Z}(U_Y \mid X_A,Z_A).
    \label{eq:scm_direct}
\end{align}
Its induced observational law is
\begin{equation}
    P_A(dz,dx,dy)
    = P_Z(dz)\,P_{X\mid Z}(dx\mid z)\,P_{Y\mid X,Z}(dy\mid x,z)
    = P_{XYZ}(dz,dx,dy),
    \label{eq:direct_law}
\end{equation}
by the chain rule for probability kernels.  When
$P_{Y\mid X,Z}$ varies with $x$, this model has a nonzero controlled
causal effect of $X$ on $Y$ conditional on $Z$.

\medskip\noindent
\textbf{World B: common-cause confounding.}
For each $z$, choose a latent variable $U_C \sim \mathrm{Unif}(0,1)$ and
measurable functions $g_X,g_Y$ such that
\begin{equation}
    (g_X(z,U_C),g_Y(z,U_C)) \sim P_{XY\mid Z=z}.
    \label{eq:common_cause_kernel}
\end{equation}
This is possible by the randomization lemma for regular conditional
distributions.  Define
\begin{align}
    Z_B &= Q_Z(U_Z), \nonumber\\
    C_B &= (Z_B,U_C), \nonumber\\
    X_B &= g_X(Z_B,U_C), \nonumber\\
    Y_B &= g_Y(Z_B,U_C).
    \label{eq:scm_confounding}
\end{align}
Then $X_B$ and $Y_B$ have no directed edge between them; their dependence
is generated by the common cause $C_B$, whose observed component is $Z_B$.
The induced law is
\begin{equation}
    P_B(dz,dx,dy)
    = P_Z(dz)\,P_{XY\mid Z}(dx,dy\mid z)
    = P_{XYZ}(dz,dx,dy).
    \label{eq:confounding_law}
\end{equation}
If the analyst additionally declares that the observed $Z$ is the complete
common cause and that the exogenous errors are independent, this construction
specializes to $Z \to X$ and $Z \to Y$ exactly when
$X \perp\!\!\!\perp Y\mid Z$.  Without that completeness declaration, the
latent component $U_C$ is observationally invisible.

\medskip\noindent
\textbf{World C: reverse causation $Y \to X$.}
Define the SCM
\begin{align}
    Z_C &= Q_Z(U_Z), \nonumber\\
    Y_C &= Q_{Y\mid Z}(U_Y \mid Z_C), \nonumber\\
    X_C &= Q_{X\mid Y,Z}(U_X \mid Y_C,Z_C).
    \label{eq:scm_reverse}
\end{align}
Again, by the chain rule,
\begin{equation}
    P_C(dz,dx,dy)
    = P_Z(dz)\,P_{Y\mid Z}(dy\mid z)\,P_{X\mid Y,Z}(dx\mid y,z)
    = P_{XYZ}(dz,dx,dy).
    \label{eq:reverse_law}
\end{equation}

Equations~\eqref{eq:direct_law}, \eqref{eq:confounding_law}, and
\eqref{eq:reverse_law} show that the three SCMs induce identical
observational distributions over $(X,Y,Z)$ while assigning different
counterfactual meanings to the same association.  Let
$T_n = T((X_i,Y_i,Z_i)_{i=1}^n)$ be any statistic or verification
procedure based only on observational samples.  Since the sample law is
$P_{XYZ}^{\otimes n}$ in all three worlds, the distribution of $T_n$ is
identical in all three worlds.  Therefore $T_n$ cannot verify that World A,
World B, or World C is the true counterfactual data-generating process.
The missing information is precisely the unobserved counterfactual
response under interventions.
\end{proof}

\begin{corollary}[No-Assumptions Causal Claims Have Zero Audit Content]
\label{cor:no_assumption_causal_claim}
Let $\mathcal{A}_{\mathrm{id}}$ be the set of declared identification
strategies, including randomized assignment, a valid backdoor adjustment
set, a valid front-door mediator, a valid instrumental variable, regression
discontinuity, or another explicit design that identifies
$P(Y\mid \mathrm{do}(X))$.  A causal claim about $X$ and $Y$ that declares no element
of $\mathcal{A}_{\mathrm{id}}$ is epistemically equivalent, in the SCX sense,
to reporting the observational law $P_{XYZ}$ alone.
\end{corollary}

\begin{proof}
By Theorem~\ref{thm:causal_unident}, the equivalence class of SCMs compatible
with $P_{XYZ}$ contains at least one direct-causal, one common-cause, and one
reverse-causal interpretation.  A declared identification strategy is exactly
an assumption that removes some members of this equivalence class.  If no such
strategy is declared, the identified set for the causal estimand contains all
counterfactual values generated by the compatible SCMs.  The claim therefore
does not shrink the equivalence class beyond what was already implied by
$P_{XYZ}$.  This is the causal instance of Theorem~\ref{thm:lucas_unident}:
without declared invariance or identification assumptions, the source of the
observed association is unidentifiable.
\end{proof}

\end{document}
